\chapter{Зарождение: кто начинает первым}\label{ch:nucleation}

\section{Классическая теория зарождения}

Новая фаза появляется не сразу целиком, а маленьким зародышем. У зародыша есть два вклада в
энергию (рис.~\ref{fig:F13_CNT}): выигрыш объёма $-\Delta\,V(r)$ ($\Delta$ --- движущая сила на
единицу объёма из гл.~\ref{ch:solubility}) и проигрыш поверхности $+\gamma\,S(r)$ ($\gamma$ ---
энергия межфазной границы). Объём растёт как $r^3$, поверхность --- как $r^2$, поэтому у малых
зародышей побеждает поверхность, у больших --- объём. Между ними --- горб, \emph{барьер}.

\fig{F13_CNT}{Энергия зародыша от его размера при трёх движущих силах. Чем больше $\Delta$, тем
ниже и левее барьер. Скорость зарождения экспоненциально зависит от высоты барьера.}

Для сферы $\Delta G(r)=4\pi r^2\gamma-\frac43\pi r^3\Delta$, максимум при $r^*=2\gamma/\Delta$:
\begin{equation}
\Delta G^*=\frac{16\pi\gamma^3}{3\Delta^2},\qquad J=\nu\,\exp\!\left(-\frac{\Delta G^*}{kT}\right).
\end{equation}
Для пластинки или линзы меняется только числовой множитель («фактор формы»), а главное ---
$\Delta G^*\propto\gamma^3/\Delta^2$ --- остаётся.

\begin{exercise}
Выведите $r^*$ и $\Delta G^*$ для сферы. Во сколько раз упадёт барьер, если движущая сила
вырастет на 10\%?
\end{exercise}

\section{Как теория записана в модели}

Мы не знаем $\gamma$, $\nu$ и форму зародыша гидрида достаточно точно, чтобы считать $J$ из
первых принципов. Зато знаем, где зарождение становится заметным: на линии TSSP. Поэтому
барьер записан относительно неё. Пусть на TSSP движущая сила равна $\Delta_0$, а барьер ---
$B\,kT$. Так как барьер $\propto1/\Delta^2$, при другой движущей силе он равен $B\,kT\,(\Delta_0/\Delta)^2$:

\begin{keyf}[Скорость зарождения в клетке]
\[
r=\nu_c\exp\!\left[-B\left(\frac{\Delta_0^2}{\Delta^2}-1\right)\right],\qquad
\Delta=\Delta_0+n_HRT\ln\frac{c}{c_{\TSSP}}+\eps_n g+n_H\frac{Q}{T}\,\delta T .
\]
$B=80$, $\Delta_0=15$~МПа, $\nu_c=10^{-6}$~с$^{-1}$ на клетку (на TSSP $r=\nu_c$).
\end{keyf}

Четыре слагаемых $\Delta$ --- четыре физики модели:
\begin{enumerate}
\item $\Delta_0$ --- запас движущей силы на линии TSSP;
\item $n_HRT\ln(c/c_{\TSSP})$ --- химия: насколько раствор пересыщен в этой клетке сейчас
  (водород диффундирует и расходуется, гл.~\ref{ch:hydrogen});
\item $\eps_n g$ --- поле соседних пластинок (гл.~\ref{ch:misfit}--\ref{ch:plasticity});
\item $n_H\frac{Q}{T}\,\delta T$ --- сдвиги температуры выпадения в градусах: от нагрузки, от
  ориентации зерна, от границы (ниже).
\end{enumerate}

\fig{F14_Rate}{Логарифм скорости зарождения от движущей силы. У линии TSSP ($\Delta=\Delta_0$)
кривая очень крутая: $+1$~МПа --- скорость больше в $e^{11}$ раз, $+1\,^\circ$C --- в $e^{58}$.}

\begin{idea}
Зарождение --- не «выгодно/невыгодно», а лавина. Крутизна у TSSP $d\ln r/d\Delta=2B/\Delta_0\approx11$
на мегапаскаль, то есть $\approx58$ на градус. Семейство зёрен, которое на долю градуса раньше
достигло своей температуры выпадения, зарождается в миллиарды раз чаще и забирает водород.
Отсюда резкий порог.
\end{idea}

\begin{exercise}
Проверьте $d\ln r/d\Delta=2B/\Delta_0$ при $\Delta=\Delta_0$. Во сколько раз отличаются скорости
зарождения двух зёрен с разницей температуры выпадения 0.2\,$^\circ$C?
\end{exercise}

\section{Нагрузка и ориентация зерна --- в градусах}

\paragraph{Нагрузка.} В кинетической модели влияние приложенного напряжения на выбор ориентации
взято из опыта, а не из упругой формулы (гл.~\ref{ch:misfit}): синхротронные измерения Vizcaíno
и др. (2014) показали, что в зёрнах с осью $c$ вдоль нагрузки гидрид выпадает раньше на
$0.08\pm0.02\,^\circ$C на мегапаскаль нормального к пластинке напряжения. Для пластинки под углом
$\psi$ к окружности нормальное напряжение $\sigma_{nn}=\sigma_\theta\sin^2\psi$, и
\begin{equation}
\delta T_{\text{нагр}}=0.08\,^\circ\text{C/МПа}\cdot\sigma_\theta\sin^2\psi .
\end{equation}

\paragraph{Фора зёрен.} Даже без нагрузки зёрна разной ориентации выпадают при разных температурах:
у Vizcaíno разница между семействами 5\,$^\circ$C (Zr-2.5Nb). В модели это \emph{фора}
$\delta T_{\text{фора}}=b\cos^2\psi$ зёрнам, где пластинка окружная (ось $c$ по радиусу).
У Zr--Nb $b=5\,^\circ$C (измерено), у холоднодеформированного Zircaloy-4 --- 10.9--12.4\,$^\circ$C
(подобрано по порогу Cinbiz). Природа форы, по книге Пульса, --- скорее остаточные напряжения между
зёрнами после прокатки ($\sim100$--$150$~МПа на базисных плоскостях), чем плотность дислокаций:
холодная деформация сама TSSP не сдвигает.

\section{Границы зёрен}

\fig{F15_GB}{Зародыш на границе --- линза. Граница сама стоит энергии $\gamma_{\alpha\alpha}$, и
зародыш, «съедая» её кусок, экономит; барьер умножается на фактор формы $f(\theta)<1$.}

Опыт (Son и др. 2026, EBSD) показывает, что в холоднодеформированном металле большая часть
гидридов лежит на границах зёрен, причём гидрид держит ориентационное соотношение с одним из
соседей, а граница выбирается близкой к базисной плоскости этого соседа (Qin 2011). В модели:
\begin{itemize}
\item на гранях Вороного зарождение разрешено, если след грани отличается от базисного следа
  одного из соседей не больше чем на $15^\circ$;
\item барьер на границе ниже: зародыш-линза с углом смачивания
  $\cos\theta=\gamma_{\alpha\alpha}/(2\gamma_{\alpha\delta})$ имеет фактор формы
  \begin{equation}
  f(\theta)=\frac{(2+\cos\theta)(1-\cos\theta)^2}{4};
  \end{equation}
  с энергиями Jo и др. (2025) $\gamma_{\alpha\alpha}=0.14$, $\gamma_{\alpha\delta}=0.18$~Дж/м$^2$
  получаем $\cos\theta=0.39$, $f=0.22$;
\item барьер $fB$ при той же нормировке даёт заметное зарождение уже при $\Delta=\Delta_0\sqrt f$,
  на $\Delta_0(1-\sqrt f)\approx8$~МПа раньше, то есть на $\approx1.5\,^\circ$C. Это и есть фора
  границы $\delta T_{\text{гр}}$, никаких подгонок.
\end{itemize}

\begin{warnbox}
Меньший барьер на границе заменён сдвигом начала зарождения на 1.5\,$^\circ$C при прежней крутизне.
Честный барьер $fB$ дал бы и другую крутизну. Плотность мест на границах (их меньше, чем клеток в
теле зерна) учтена только тем, что граница --- полоса в одну клетку; поэтому доля межзёренных
гидридов зависит от шага сетки и её надо считать оценкой.
\end{warnbox}

\begin{exercise}
Посчитайте $f$ для $\cos\theta=0.39$ и сдвиг $\Delta_0(1-\sqrt f)$ в МПа и в градусах.
Как изменится фора границы, если $\gamma_{\alpha\alpha}$ на 30\% больше?
\end{exercise}

\section{Гонка семейств и порог}

Теперь всё складывается (рис.~\ref{fig:F16_Race}). Окружным пластинкам фора даётся всегда
(10.9\,$^\circ$C в теле зерна, ещё 1.5\,$^\circ$C на подходящих границах), радиальным --- только
нагрузкой, $0.08\,\sigma_\theta$. Пока нагрузка мала, окружные выпадают первыми и забирают
водород, который за минуты стекается к ним со всего поля (гл.~\ref{ch:hydrogen}). Когда
$0.08\,\sigma_\theta$ превысит фору, первыми начинают радиальные.

\fig{F16_Race}{Фора по температуре выпадения против окружного напряжения. Пересечения ---
пороги переориентации: 136~МПа, когда радиальные обгоняют окружные в теле зерна, и 155~МПа ---
окружные на границах.}

\begin{keyf}[Порог]
\[
\sigma_{\text{порог}}\approx\frac{\text{фора окружного семейства}}{0.08\,^\circ\text{C/МПа}} .
\]
\end{keyf}

\begin{exercise}
Посчитайте порог для форы 12.4\,$^\circ$C (без границ) и для 5\,$^\circ$C + 1.5\,$^\circ$C (Zr--Nb
с границами). Сравните с опытом: Cinbiz $\approx155$~МПа (Zircaloy-4), Son: начало около 80~МПа (Zr--Nb).
\end{exercise}

\section{Случайность: сколько зародышей за шаг}

Скорость $r(\mathbf x)$ задана в каждой клетке. За шаг $\delta t$ ожидаемое число зародышей
$\lambda=\sum_{\mathbf x}r(\mathbf x)\,\delta t$. Зарождение --- редкие независимые события, поэтому
их число разыгрывается по Пуассону со средним $\lambda$, а места --- с вероятностями,
пропорциональными $r(\mathbf x)$. Если $\lambda$ больше 6, шаг по времени дробится (не больше чем
в 20 раз): зародыши одного шага друг друга «не видят», и их не должно быть слишком много.

\begin{exercise}
Почему зародыши одного шага не должны видеть друг друга, и что случилось бы, если бы шаг был
настолько большим, что на нём рождалась бы сотня пластинок? (Подсказка: поле и водород
пересчитываются только между шагами.)
\end{exercise}

\begin{codebox}
\textbf{Где в коде.} \code{ca\_kinetic.py}: \code{KParams} (\code{B}, \code{Delta0}, \code{nu\_c},
\code{app\_dT}, \code{bias\_dT}, \code{dT\_s}, \code{gb}, \code{gb\_dT}, \code{gb\_tol}, \code{lam\_max});
движущая сила и скорость --- главный цикл \code{run\_kinetic}; грани зёрен ---
\code{ca\_hydride.make\_grains(gb=True)}.
\end{codebox}
